Esta seção reúne os conceitos em que a comparação se apoia, do problema ao
vocabulário com que ela se faz.
Começa pelo que os eventos e o estado pedem do código de uma interface e
pela diferença entre o imperativo e o declarativo.
Depois trata do modelo de programação, que é a unidade comparada, e da
\gls{pf} e da \gls{pr}, de que vêm os modos de cada eixo.
Termina pelas \glspl{dc}, com que a análise descreve a diferença, e pelos
trabalhos relacionados.
\section{Eventos, estado e \emph{callbacks}}
\label{sec:eventos-estado-e-callbacks}

Um programa orientado a eventos reage ao que acontece no ambiente externo: o
clique num botão, a tecla digitada, a mensagem do servidor
\cite[p. 86:1]{madsen2017}.
Os eventos entram numa fila, e o programa os trata um de cada vez
\cite[p. 86:1]{madsen2017}, sem garantia da sequência em que chegam:
como a interação depende de eventos externos, que podem ocorrer em
qualquer ordem, entender o programa exige seguir várias cadeias de causa
e efeito \cite[p. 125]{martin2022}.
Cada evento pode alterar o \emph{estado} da aplicação.
Uma barra de ferramentas, por exemplo, fica encaixada acima da área de
trabalho até que o usuário a segure pelo topo, gesto que muda o estado
dela: a barra se solta e passa a seguir o mouse \cite[p. 230]{oney2012}.

Nem todo evento vem do usuário: a chegada de um pacote da rede e o tique do
relógio também são eventos, e cada evento faz rodar a parte do programa que
responde a ele \cite[p. 8]{krishnamurthi2019}.
O desafio do programador é escrever, no código, como o estado passa de um
evento a outro \cite[p. 8]{krishnamurthi2019}.

O programa também precisa ser robusto: tratar qualquer entrada sem parar
de funcionar e, como o usuário espera, permitir que ele cancele e desfaça
a maioria das ações \cite[p. 75, 79 e 80]{myers1994}.

O \emph{callback} é síncrono quando a função que o recebe o chama antes de
retornar, e assíncrono quando sua execução fica para algum tempo depois de a
função retornar \cite[p. 2]{gallaba2015}.
O navegador executa a aplicação numa \emph{thread} só, e, para que ela continue
respondendo ao usuário, os programadores recorrem às operações assíncronas
\cite[p. 368-369]{alimadadi2014}.
O \emph{callback} de um clique é desse tipo: registrado quando a tela é montada,
só roda quando o usuário clica, e os eventos do \gls{dom} estão entre as
operações assíncronas do navegador \cite[p. 5]{gallaba2015}.

Em JavaScript, os \emph{callbacks} são frequentes: num estudo de 138 programas, em
média uma em cada dez definições de função recebe um \emph{callback}, e a maioria dos
\emph{callbacks} é aninhada \cite[p. 1]{gallaba2015}.
No código do navegador, em média 9\% das chamadas de função passam um
\emph{callback}, e 72\% desses \emph{callbacks} são assíncronos \cite[p. 6]{gallaba2015}.

O \emph{callback} fragmenta o fluxo de controle: a sequência de passos que o
programador tem em mente para uma operação fica dividida entre vários
\emph{callbacks} \cite[p. 134-135]{fischer2007}.
Numa rede de \emph{callbacks}, a ordem em que os eventos chegam pode depender da
ordem em que os \emph{callbacks} foram registrados, o que a torna difícil de prever
\cite[p. 8]{blackheath2016}.
Em 2026, o tratamento de eventos nas interfaces ainda se apoia em
\emph{callbacks} ou \emph{listeners}, que tornam difíceis de exprimir e de manter os
fluxos de interação complexos \cite[p. 8]{grolaux2026}.

O JavaScript ganhou alternativas ao \emph{callback}: as \emph{promises}, que encadeiam
computações assíncronas \cite[p. 86:1]{madsen2017}, e as palavras-chave
\emph{async} e \emph{await}, que permitem escrever operações assíncronas em estilo
sequencial \cite[p. 9]{grolaux2026}.
As \emph{promises} achatam o aninhamento dos \emph{callbacks}, mas a lógica complexa
continua presa em cadeias de \emph{callbacks} \cite[p. 9]{grolaux2026}.
As \emph{promises} têm semântica complexa, e o fluxo de controle e de dados do
código que as usa dificulta a compreensão e leva facilmente a erros
\cite[p. 162:1]{alimadadi2018}.
Numa implementação em JavaScript que guarda o estado da aplicação em variáveis
globais, o \emph{async} e o \emph{await} não ajudam a coordenar esse estado
\cite[p. 535 e 538]{berry2020}.

\emph{Threads} exprimem sequência, e eventos exprimem dependência, e muitos
problemas vêm de escrever uma com o recurso da outra
\cite[p. 16]{blackheath2016}.
Os dois modos imperativos da orquestração correspondem a esses dois
recursos.
Na orquestração por atribuição, cada \emph{callback} trata um evento e grava o
valor no estado; na sequencial, o código espera o próximo evento e segue
em ordem, como permitem o \emph{async} e o \emph{await}.
\section{Imperativo e declarativo}
\label{sec:imperativo-e-declarativo}

Nas linguagens imperativas, o estado é implícito e muda por comandos,
como a atribuição de um valor a uma variável, e essas alterações são os
\emph{efeitos colaterais} dos comandos \cite[p. 361]{hudak1989}.
Por isso, para controlar o estado com precisão e de forma determinística,
elas costumam fixar a ordem em que os comandos se executam
\cite[p. 361]{hudak1989}.
Numa interface, é o programador quem escreve essa ordem: para cada alteração
de estado, ele precisa escrever a atualização de tudo o que depende do
valor alterado, e esse código é complexo e propenso a erros, como o que
muda o estado na hora errada ou na ordem errada
\cite[p. 52:2]{bainomugisha2013}.
Para \textcite[p. 929]{edwards2009}, boa parte da sequência que o
programador escreve poderia ficar a cargo da execução, ao menos nas
aplicações interativas.

O código sequencial tem o problema oposto ao do \emph{callback}: a linguagem
obriga a fixar a ordem de execução mesmo quando ela não importa, e quem lê
precisa supor que importa \cite[p. 8-9]{moseley2006}.
Para \textcite[p. 1]{moseley2006}, essa preocupação com o fluxo de controle
está, ao lado do volume de código e do tratamento do estado, entre as
principais causas da complexidade dos sistemas de grande porte.
Os autores distinguem a complexidade essencial, própria do problema, da
acidental, todo o resto, e propõem reduzir a acidental com a \gls{pf} e
o modelo relacional de dados \cite[p. 1 e 21]{moseley2006}.

Nas linguagens declarativas, não há estado implícito, e o programa se
escreve com expressões \cite[p. 361]{hudak1989}.
A programação declarativa costuma ser descrita como a que expressa “o
quê” e não “o como” da computação, embora a diferença seja de grau
\cite[p. 361]{hudak1989}.

O declarativo não é uma propriedade absoluta, e sim uma questão de grau
\cite[p. 406]{vanroy2004}.
Para dizer o grau de um código, este trabalho segue o critério das
\glspl{dc}, o da informação que a notação deixa à vista.
As linguagens procedurais facilitam extrair a ordem dos passos, e as
declarativas, as condições em que cada coisa acontece
\cite[p. 3]{green1989}; uma linguagem com GOTO destaca a ordem de
execução dos comandos à custa da informação declarativa
\cite[p. 136]{green1996d}.
Em cada eixo, o modo imperativo deixa à vista os passos, e o declarativo,
a relação.
Na propagação, a manual mostra os passos que atualizam a tela, e as
declarativas, como a tela depende do estado.
Na orquestração, a atribuição e a sequencial mostram os passos que alteram o estado, e a
mensagem e a acumulação, o estado seguinte como função do evento.

O grau também se dá a cada parte do código, e não à tecnologia inteira.
Numa classificação das tecnologias de interface em seis categorias, até a
das declarativas escreve de forma declarativa a estrutura e a ligação com
os dados, e de forma imperativa o comportamento, com tratadores de eventos
\cite[p. 765]{mijailovic2014}.
Mesmo assim, as fontes costumam dar o rótulo à tecnologia inteira: para
\textcite[p. 93]{disch2025}, a maioria dos \emph{frameworks} de interface
atuais é imperativa, porque o estado é mutável e compartilhado, e os
\emph{callbacks} o leem e o alteram.
Por isso, este trabalho dá o grau a cada eixo do modelo de programação, e
não à tecnologia.
\section{Paradigma e modelo de programação}
\label{sec:paradigma-e-modelo-de-programacao}

O modelo de programação não é o paradigma, noção imprecisa que o modelo
de computação torna precisa \cite[p. xiii]{vanroy2004}.
Também não é a biblioteca nem o \emph{framework}, que classificam o produto: o
Web Component, recurso do próprio navegador \cite{mozilla2026}, não é
nenhum dos dois e tem um modelo de programação como as outras
tecnologias.

Um \emph{paradigma de programação} é uma abordagem da programação baseada numa
teoria matemática ou num conjunto coerente de princípios
\cite[p. 10]{vanroy2009}.
Cada modelo de programação aplica conceitos de um ou mais paradigmas:
em quase todo problema, cada parte pede o seu conjunto de conceitos
\cite[p. 10]{vanroy2009}, e as implementações deste trabalho também os
misturam.

Na classificação de \textcite[p. 14]{vanroy2009}, uma linguagem só
entra num paradigma quando a linguagem núcleo dela o suporta, e não basta
que haja bibliotecas escritas nela para isso.
Os modelos de programação deste trabalho se escrevem todos na mesma
linguagem, o TypeScript, e só as bibliotecas e a própria plataforma os
separam.
Por isso a unidade comparada é o modelo de programação, e não o
paradigma.
Entender os conceitos certos melhora a programação mesmo numa linguagem
que não os suporta, como a programação orientada a objetos é possível em
C, com a atitude certa do programador \cite[p. 10]{vanroy2009}.

A própria noção de paradigma é contestada.
Para \textcite[p. 1]{krishnamurthi2019}, a literatura se perde num
conceito mal definido e até confuso.
As linguagens não se organizam em taxonomias, porque as mais usadas
misturam conceitos de vários paradigmas \cite[p. 7]{krishnamurthi2019}.
Ser imperativo também não impede ser reativo: a reatividade diz como o programa se comporta quando um valor se
altera, e não como ele se organiza \cite[p. 7]{krishnamurthi2019}.
Na comparação de \textcite[p. 23 e 56]{kiss2014} entre a \gls{poo} e
a \gls{pf} em sete interfaces, o \emph{toolkit} teve o papel principal, e as
diferenças de paradigma e de linguagem pesaram menos.

O mesmo texto de programa também pode ter outra semântica.
A ordem que ele fixa está na semântica da linguagem, e não no texto, que
seria um
programa válido numa linguagem que não executasse os comandos na ordem
em que se escrevem \cite[p. 9]{moseley2006}.
É o caso do JSX, que se escreve com a mesma sintaxe no React e no Solid: o
React reexecuta o componente quando o estado se altera, e o Solid
atualiza só os trechos da tela que dependem dele
\cites[][ p. 289:38]{lee2025}{solidjs2026}.
Por isso, a mesma montagem da tela pode pertencer a dois modelos de
programação.
\section{Programação funcional}
\label{sec:programacao-funcional}

A \gls{pf} se chama assim porque o programa consiste inteiramente de
funções \cite[p. 1]{hughes1990}.
São funções matemáticas: chamadas com os mesmos argumentos, dão sempre os
mesmos resultados \cite[p. 26]{vanroy2009}.
O componente do React segue essa ideia: é uma função que, idealmente,
chamada com os mesmos argumentos, exibe a mesma tela
\cite[p. 289:1]{lee2025}.

Uma função também pode receber outra como argumento.
A função que recebe é de ordem superior, e a programação de ordem
superior é o conjunto das técnicas que essa possibilidade abre
\cite[p. 177]{vanroy2004}.
O \emph{callback} é uma delas.
Com as listas, a ordem superior toma o lugar do laço: o \texttt{map} calcula uma
lista nova aplicando uma função a cada elemento da antiga, e o \texttt{filter}
guarda os elementos para os quais uma função dá verdadeiro
\cite[p. 190-191]{vanroy2004}.
O \emph{fold} combina os elementos num valor só: a soma de uma lista se escreve
como um padrão geral de recursão, o \emph{fold}, mais a parte própria do
problema, a função que soma dois números \cite[p. 5]{hughes1990}.

O laço \texttt{for} diz “o como” de uma transformação, e o \texttt{map} e o \texttt{filter}
dizem “o quê”.
O Código \ref{code:precos-for} formata os preços positivos de uma lista
com o laço, que percorre os índices, testa cada preço e acrescenta o
formatado a uma lista que começa vazia.
O Código \ref{code:precos-map} faz o mesmo com o \texttt{filter} e o \texttt{map}, que
recebem a função a aplicar, e a lista nova se define pelo que ela contém.
É um passo de um modo a outro: o resultado não muda, e muda o que o código
deixa à vista.

  \begin{listing}[htbp]
    \caption{Os preços positivos formatados com o laço \texttt{for}.}
    \begin{minted}{typescript}
function formatarPreco(preco: number): string {
  return "R$ " + preco.toFixed(2).replace(".", ",");
}

const precos = [50.45, 0, 20.99, 3.44];
const formatados: string[] = [];
for (let i = 0; i < precos.length; i++) {
  if (precos[i] > 0) {
    formatados.push(formatarPreco(precos[i]));
  }
}
    \end{minted}
    \small Fonte: elaborado pelo autor.
    \label{code:precos-for}
  \end{listing}

  \begin{listing}[htbp]
    \caption{Os mesmos preços formatados com o \texttt{filter} e o
      \texttt{map}, com a \texttt{formatarPreco} do
      Código~\ref{code:precos-for}.}
    \begin{minted}{typescript}
const precos = [50.45, 0, 20.99, 3.44];
const formatados = precos
  .filter((preco) => preco > 0)
  .map(formatarPreco);
    \end{minted}
    \small Fonte: elaborado pelo autor.
    \label{code:precos-map}
  \end{listing}
\section{Programação reativa}
\label{sec:programacao-reativa}

O \emph{Observer Pattern} separa os observadores, que consomem os eventos, dos
observáveis, que os produzem \cite[p. 953]{salvaneschi2015}.

A \gls{pr} oferece abstrações para os valores que variam no tempo e
cuida sozinha das dependências entre eles
\cite[p. 52:1]{bainomugisha2013}.
\textcite[p. 52:2]{bainomugisha2013} chamam as planilhas de “possivelmente a
linguagem de programação mais utilizada por usuários finais” (tradução nossa).
A maior parte da pesquisa em \gls{pr}, dizem ainda, descende do Fran, uma
linguagem da \gls{pf} do fim dos anos 1990.
A \gls{pr} traz duas abstrações: os \emph{behaviours}, para os valores que variam de forma contínua
no tempo, e os eventos, para os valores discretos
\cite[p. 52:2]{bainomugisha2013}.

Na soma \texttt{var3 = var1 + var2}, a programação sequencial imperativa guarda
em \texttt{var3} a soma dos valores iniciais, mesmo que \texttt{var1} ou \texttt{var2} recebam
depois outro valor, a não ser que o programador escreva de novo a
atribuição \cite[p. 52:3]{bainomugisha2013}.
Na \gls{pr}, \texttt{var3} se recalcula sempre que \texttt{var1} ou \texttt{var2} se
alteram \cite[p. 52:3]{bainomugisha2013}.
É o que a planilha faz com as fórmulas.

Na \gls{pr}, a ordem das atualizações não se escreve: o ambiente de
execução propaga a alteração de um valor às computações que dependem dele
e as recalcula na ordem do grafo de dependências, cada uma depois das que
ela usa \cite[p. 52:5 e 52:7]{bainomugisha2013}.
As fontes descrevem também como essa propagação se executa.
No modelo \emph{pull}, a computação que precisa do valor o busca na fonte, e no \emph{push},
a fonte empurra o valor novo às computações que dependem dele; há
linguagens que combinam os dois \cite[p. 52:5-52:6]{bainomugisha2013}.
Este trabalho não classifica os modos por essa execução, e sim pelo que o
código escreve.

Nas bibliotecas de \gls{pr} para linguagens de uso geral, as duas
abstrações aparecem como \emph{signals} e fluxos de eventos.
O \emph{signal} é um conceito de linguagem para escrever, de forma declarativa,
as dependências funcionais entre valores \cite[p. 1126]{salvaneschi2017},
e é dele que vem a propagação granular.
Os fluxos de eventos vêm com operadores que os transformam, filtram e
combinam \cite[p. 1126]{salvaneschi2017}, e deles vêm a propagação por
combinadores e a orquestração por acumulação, em que o estado é o \emph{fold}
dos eventos.

As aplicações orientadas a eventos e interativas, na literatura de
\gls{pr}, também se chamam reativas \cite[p. 1125]{salvaneschi2017}, mas
a palavra tem sentidos diferentes em linguagens de programação e em sistemas
distribuídos \cite[p. 5:1]{myter2019} e já nomeia tecnologias como o React
\cite[p. 13]{grolaux2026}.
Para \textcite{berry1989}, reativo é o programa que mantém a interação
com o ambiente no ritmo do ambiente, e não no próprio, como os que tratam
a interface entre a pessoa e a máquina.
\textcite[p. 564]{salvaneschi2014} põem o React entre as bibliotecas que
implementam princípios da \gls{pr}, mas nele o código não escreve a
dependência entre os valores: quando o estado se altera, o componente
inteiro se reexecuta.
Por isso, este trabalho chama essas aplicações de interativas e reserva
“reativo” para a \gls{pr}.

Um contador mostra os três modos da orquestração que o núcleo percorre,
com a tela escrita do mesmo jeito nos três, para só a orquestração mudar.
Nos três, \texttt{botao} e \texttt{caixa} são o botão e o elemento que mostra o valor.
No Código \ref{code:contador-atribuicao}, por atribuição, o \emph{callback}
do clique soma um à variável.
No Código \ref{code:contador-mensagem}, por mensagem, o clique vira uma
mensagem que uma função pura, \texttt{atualizar}, transforma no valor seguinte,
como o \emph{reducer} que o \texttt{useReducer} do React recebe \cite{meta2026d}.
No Código \ref{code:contador-acumulacao}, por acumulação, o valor se
define sobre o fluxo dos cliques, com o \texttt{scan} do RxJS, uma variação do
\emph{fold} \cite[p. 1530]{zimmerle2025}, e nenhuma variável o guarda.

  \begin{listing}[htbp]
    \caption{O contador com a orquestração por atribuição.}
    \begin{minted}{typescript}
let valor = 0;
botao.addEventListener("click", () => {
  valor = valor + 1;
  caixa.textContent = String(valor);
});
    \end{minted}
    \small Fonte: elaborado pelo autor.
    \label{code:contador-atribuicao}
  \end{listing}

  \begin{listing}[htbp]
    \caption{O contador com a orquestração por mensagem.}
    \begin{minted}{typescript}
type Mensagem = { tipo: "incrementar" };

function atualizar(valor: number, mensagem: Mensagem): number {
  switch (mensagem.tipo) {
    case "incrementar":
      return valor + 1;
  }
}

let valor = 0;
botao.addEventListener("click", () => {
  valor = atualizar(valor, { tipo: "incrementar" });
  caixa.textContent = String(valor);
});
    \end{minted}
    \small Fonte: elaborado pelo autor.
    \label{code:contador-mensagem}
  \end{listing}

  \begin{listing}[htbp]
    \caption{O contador com a orquestração por acumulação, no RxJS.}
    \begin{minted}{typescript}
import { fromEvent, scan } from "rxjs";

fromEvent(botao, "click")
  .pipe(scan((valor) => valor + 1, 0))
  .subscribe((valor) => {
    caixa.textContent = String(valor);
  });
    \end{minted}
    \small Fonte: elaborado pelo autor.
    \label{code:contador-acumulacao}
  \end{listing}
\section{Dimensões Cognitivas}
\label{sec:dimensoes-cognitivas}

As \glspl{dc}, propostas por \textcite[p. 443]{green1989}, tratam as
linguagens de programação como estruturas de informação, ou notações, e
são um vocabulário para descrever a
usabilidade das notações e de outros artefatos de informação
\cite[p. 327]{blackwell2001}.
Elas valem também para as bibliotecas de uma mesma linguagem:
\textcite[p. 1507]{zimmerle2025} avaliam por elas, com usuários, duas
bibliotecas de \gls{pr} para JavaScript.
Na usabilidade de \glspl{api}, as \glspl{dc} são um método corrente, mas o
conjunto original é tido como complexo, caro e voltado a quem projeta, e
por isso surgiram versões centradas em quem programa
\cite[p. 37-38]{lima2024}.

Para as \glspl{dc}, o que conta na notação é a estrutura, e não a
aparência.
\textcite[p. 333]{green2006} querem um método que não mude quando um
tocador de MP3 troca de pele, porque o funcionamento do dispositivo
continua o mesmo, e por isso a aparência fica fora da notação neste
trabalho.
Já a abstração muda a notação: com ela vem uma camada nova de notação,
em que as abstrações se escrevem, como o \gls{css}
\cite[p. 342]{green2006}.

A avaliação pelas \glspl{dc} é analítica, feita por especialista, sem a
participação de usuários \cite{blandford2008}, e
\textcite[p. 8]{kiss2014} também chama de analítica a abordagem por
elas.
As \glspl{dc}, porém, não são um método analítico, mas ferramentas de
discussão, um vocabulário para quem projeta ou avalia, de traço grosso,
que não dá previsões detalhadas \cite[p. 106]{blackwell2003}.
A avaliação analítica serve de alternativa aos estudos com usuários, que
comparam linguagens a custo alto e com resultados em geral
inconclusivos \cite[p. 9]{sadowski2011}, e as \glspl{dc} tomam
menos tempo do projetista que esses estudos \cite{dagit2006}.
A seção \ref{sec:as-dimensoes} define as 14 dimensões, e a seção
\ref{sec:a-avaliacao} diz como este trabalho as aplica.
\section{Trabalhos relacionados}
\label{sec:trabalhos-relacionados}

Três trabalhos compararam formas de programar pelas \glspl{dc}, ou por
questionários derivados delas, e dois avaliaram por elas uma biblioteca
de controles de interface.
\textcite{mernik2009} compararam pelas \glspl{dc}, com as respostas de 36
estudantes de graduação a perguntas sobre os dois, uma linguagem de
domínio específico para descrever interfaces gráficas, o \gls{xaml}, e
uma biblioteca de aplicação, o Windows Forms em C\#.
\textcite[p. v, 11 e 104]{kiss2014} comparou a \gls{poo} e a
\gls{pf} em sete interfaces gráficas, com JavaFX e ScalaFX, e em parte
delas com Scala.Rx, ReactFX e Elm, e deixou como trabalho futuro
comparar os \emph{frameworks} de interface do navegador.
O Elm, que roda no navegador, ficou em quatro das sete interfaces, porque
faltou tempo para as outras \cite[p. 96]{kiss2014}.
Com o ReactFX, e mais ainda com o Elm, o nível de abstração ficou mais
exigente, e a difusão e a viscosidade, um pouco maiores; as dependências
ocultas e a propensão a erros, por outro lado, diminuíram
\cite[p. 97]{kiss2014}.
É um resultado anterior com que a análise deste trabalho pode se
comparar.

\textcite[p. 1510 e 1513]{zimmerle2025} avaliaram com usuários duas
bibliotecas de \gls{pr} para JavaScript, RxJS e Bacon.js, em tarefas de
requisições \gls{http}, sem interface gráfica, por um questionário
que mede cinco dimensões tiradas das \glspl{dc}
\cite[p. 1507]{zimmerle2025}.
\textcite{narechania2025} e \textcite{verma2026} avaliam pelas
\glspl{dc} a notação da própria biblioteca de controles, feita em
Angular e em React, numa autoavaliação dos autores.
Não avaliam a notação do Angular nem a do React: do React, a
autoavaliação só diz que a notação da biblioteca é consistente com ele
\cite{verma2026}.
Fora das \glspl{dc}, \textcite[p. 15]{grolaux2026} compararam
conceitualmente, numa prova de conceito em JavaScript, o \emph{async} e o
\emph{await} com os \emph{listeners}, o barramento de eventos, os fluxos
funcionais e a \gls{pr} no tratamento dos eventos de interface.

\textcite[p. 1125-1134]{salvaneschi2017} testaram uma afirmação repetida,
ainda sem evidência empírica suficiente, de que a \gls{pr} torna os
programas mais fáceis de compreender que o \emph{Observer Pattern}, o padrão
tradicional da \gls{poo} para aplicações interativas.
No experimento controlado, 127 estudantes do quarto ano de Ciência da
Computação foram divididos ao acaso em dois grupos.
Um leu dez programas em Scala escritos com \gls{pr}, e o outro, os
mesmos programas com o \emph{Observer Pattern}, e cada participante respondia
a uma pergunta de múltipla escolha sobre o comportamento de cada
programa.
Quem leu os programas com \gls{pr} acertou mais (médias de 8,13 e 7,05
acertos em dez, com diferença significativa) e respondeu mais rápido em
seis das dez tarefas, sem diferença significativa nas outras quatro.
Só no grupo do \emph{Observer Pattern} os acertos acompanharam o nível de
habilidade de programação dos participantes.
Nenhum desses trabalhos compara, sobre as mesmas interfaces e em
tecnologias web, modelos de programação que vão do imperativo ao
declarativo nos dois eixos.
